%% kapa: exam=1 round=22 year=2011 period=1 form=A subject=economics subject_ko=경제원론 seq=1 number=41
%% source: Q-Net 감정평가사 기출문제 (cst00309) artlSeq=5036790 file=extracted/1008/5036790/2011년 제22회 감정평가사 1차시험 기출문제/1교시.hwp
%% conversion: kordoc 4.24.51 -> tools/gen1.py (auto); review=proofread-sonnet; flags=text-verified-hwp
\begin{question}{41}
X재와 Y재를 생산하는 K국가의 생산가능곡선상에는 두 개의 재화생산조합점 \\
$(x_{ 1 } ,~ y_{ 1 } )$=(200, 300)과 $(x_{ 2 } ,~ y_{ 2 } )$=(240, 290)이 있다. 다음 중 기회비용체증의 법칙이 성립하기 위한 이 생산가능곡선상의 재화생산조합점 $(x_{ 3 } ,~y_{ 3 } )$은? (단, $x_{ 1, } x_{ 2, } x_{ 3 }$는 각각 X재의 생산량, $y_{ 1, } y_{ 2, } y_{ 3 }$는 각각 Y재의 생산량)
\begin{choices}
  \item (160, 310)
  \item (160, 315)
  \item (280, 270)
  \item (280, 280)
  \item (280, 285)
\end{choices}
\end{question}
\answer{3}
